\lyseis{6.6}{Αντιστρόφως ανάλογα ποσά}
\sol{1} α) Ναι: $2\cdot12=4\cdot6=6\cdot4=12\cdot2=24$. β) Όχι, είναι ανάλογα ($y=5x$). γ) Όχι: $24,\,24,\,27$. δ) Ναι: $5\cdot8=10\cdot4=20\cdot2=40$.
\sol{2} $y=36,\,18,\,12,\,9,\,6,\,4$.
\sol{3} $y=\frac{60}{5}=12$, $y=\frac{60}{15}=4$, $y=\frac{60}{0{,}5}=120$. Και $x=\frac{60}{30}=2$.
\sol{4} $\frac{1}{5}$, $2$, $\frac{4}{3}$, $0{,}8$.
\sol{5} α) Σωστό. β) Σωστό. γ) Λάθος: το $x=0$ δεν επιτρέπεται (διαίρεση με το $0$). δ) Λάθος: βρίσκονται σε καμπύλη. ε) Λάθος: $\frac{48}{12}=4$ ημέρες.
\sol{6} $y=180,\,90,\,60,\,45,\,30,\,20,\,18,\,15,\,12,\,10$ λεπτά αντίστοιχα (γινόμενο $180$).
\sol{7} $4,\ 3,\ 2{,}5,\ 2,\ 1{,}5,\ 1$ ώρες (γινόμενο $120$).
\sol{8} Διαστάσεις $1\times48$, $2\times24$, $3\times16$, $4\times12$, $6\times8$. Περίμετροι $98$, $52$, $38$, $32$, $28$. Μικρότερη η $28$ m για το $6\times8$.
\sol{9} Δουλειά $4\cdot6=24$ «βρύση-ώρες». Με $3$ βρύσες $24:3=8$ ώρες, με $8$ βρύσες $3$ ώρες.
\sol{10} $\frac{72}{8}=9$, $\frac{72}{12}=6$, $\frac{72}{18}=4$ μπισκότα.
\sol{11} $60:20=3$ ώρες και $60:30=2$ ώρες. Για $1{,}5$ ώρα: $60:1{,}5=40$ σελίδες την ώρα.
\sol{12} $\frac{40}{8}=5$ € και $\frac{40}{10}=4$ €.
\sol{13} Με περισσότερα παιδιά η δουλειά τελειώνει \emph{νωρίτερα}, όχι αργότερα. Τα ποσά είναι αντιστρόφως ανάλογα: διπλάσια παιδιά, μισός χρόνος.
\sol{14} Δεν είναι ημιευθεία, αλλά καμπύλη. Η γραμμή πλησιάζει τους άξονες αλλά δεν τους φτάνει (δεν υπάρχει $x=0$ ή $y=0$).
\sol{15} Αν η ταχύτητα διπλασιαστεί ο χρόνος γίνεται το μισό, αν τριπλασιαστεί το $\frac{1}{3}$.
\sol{16} α) ανάλογα \quad β) αντιστρόφως ανάλογα \quad γ) τίποτα από τα δύο \quad δ) αντιστρόφως ανάλογα \quad ε) τίποτα από τα δύο ($y=x^2$).
\sol{17} $y=\frac{36}{x}=36\cdot\frac{1}{x}$, άρα $\frac{y}{1/x}=36$: ανάλογα με συντελεστή $36$.
\sol{18} Ελεύθερη απάντηση. Π.χ. $1{,}2$ km: με $5$ km/h $0{,}24$ h $\approx14$ λεπτά, με $15$ km/h $0{,}08$ h $\approx5$ λεπτά.
\sol{19} Ελεύθερη απάντηση. Π.χ. ταχύτητα και χρόνος, αριθμός εργατών και ημέρες εργασίας.
